\cotefile{57}
\scanpage{287}
\makeatletter\tagsleft@true\makeatother
Calcul du \unsure{fsct|fct} $\underline{\Omega}^{1}_{X/S})$. On aura

\diagram{\begin{tikzcd}[ampersand replacement=\&,column sep=small]
X\arrow[d,"p"'] \& A\times_S B\arrow[l,"\Gamma"']\arrow[dl,"q"]\\
S \&
\end{tikzcd}}

\[
p_*(\underline{\Omega}^{1}_{X/S})
=q_*(\underline{\Omega}^{1}_{A\times_S B})^{\Gamma}
=\pi_{A*}(\underline{\Omega}^{1}_{A/S})^{\Gamma}\times\pi_{B*}(\underline{\Omega}^{1}_{B/S})^{\Gamma}
\]

Donc dans le cas \unsure{parti.} où $A$ et $B$ sont des schémas abéliens, \deleted{on trouve} $\Gamma$ opère sur $B$ via translations, on trouve

\[
\deleted{\unsure{p_*}}\qquad \omega^{1}_{X/S}\simeq\omega^{1}_{B/S}\times{\omega^{1}_{A/S}}^{\Gamma}
\tag*{a)}
\]

Si $S$ est le spectre d'un corps $k$, pour comparer $\operatorname{rang}_k H^1(X,\underline{O}_X)$ et $\operatorname{rang}_k H^0(X,\Omega^1_{X/k})$, on note \deleted{doit} que en vertu de \unsure{a)}, on aura

\[
H^0(X,\underline{\Omega}^1_X)\simeq H^0(B,\Omega^1_B)\times H^0(A,\Omega^1_A)^\Gamma
\tag*{a$'$)}
\]

\[
H^1(X,\underline{O}_X)\simeq H^1(Y,\underline{O}_Y)\times\underbrace{H^1(A,\underline{O}_A)^\Gamma}_{\mathbf{t}_{A'}^\Gamma}
\tag*{b$'$)}
\]

donc comparant avec \unsure{a)}, on aura

\[
\operatorname{rang} H^1(\unsure{X|1},\underline{O}_X)-\operatorname{rang} H^0_{\unsure{X}}(X,\Omega^1_{X/k})=\dim t_{A'}^\Gamma-\dim(t_A)'{}^\Gamma
\]

Nous voulons donner des exemples où cette \inserted{différence de} dimensions est $>0$, ou $<0$, \deleted{\unsure{enfin|et}} ou [dans le cas d'une famille] \deleted{chaque}\inserted{les} termes \deleted{\unsure{du 2ème mbr}} $\dim(t_{A'})^\Gamma$ et $\dim(t_A)'{}^\Gamma=\deleted{\unsure{coinv|\dim}}\inserted{\dim}(\mathbf{t}_A)_\Gamma$ varient d'une valeur du paramètre à l'autre \deleted{[\unsure{pour un}} [l'une spécialisant l'autre],

% archive sheet number: 286

